% REVISED: canonical-target reframe 2026-06-04
% REVISED: hostile-review action plan 2026-06-04
% REVISED: JLEO compression pass 2026-06-06 — full dictionary, fold audits,
% calibration, and the design-by-design performance table migrate to Online
% Supplement E (\ref{supp:bid_cost}); appendix keeps model summary + pointer.
%
\section{Bid-Layer Benchmark and Cost-Recall Details}
\label{app:bid_benchmark_submission}

Section~6.1 asks whether the cheap, award-layer participation signal
carries ranking information comparable to a transparent bid-moment
benchmark built from full bid microdata; this appendix gives the model
summary, and the full construction, feature dictionary, fold audits, and
design-by-design performance table are in Online Supplement~E. The
benchmark is an \emph{implemented
full-observability} diagnostic: it computes seven Imhof--Wallimann
within-tender bid-distribution moments, aggregates them to the firm, and
feeds them to a random forest evaluated by cross-validation. It is neither
a state-of-the-art horse race nor a first-stage triage tool, because the
moments require every within-tender bid to be observed and recovered.
``Cobidder'' is an adjudication-anchored exposure label throughout---an
always-loser firm that shares at least one BEC tender-item with a
BEC-active direct CADE defendant---so the benchmark \emph{ranks} that
exposure target for forensic priority and makes no claim to identify
cartelists.

%--------------------------------------------------------------------
\subsection{Data, Units, and Support}
\label{app:bid_data_submission}

The bid layer is built from one row per priced bid (firm, OC sequence,
item, \texttt{bid\_price}), retaining a bid only if
\texttt{bid\_price > 0}. The modeling unit is the \emph{firm}: per-tender
moments are shared across all firms bidding on a tender-item and then
firm-mean aggregated. Of the $16{,}843$ always-losers, $64$ are dropped
for having no priced bids or fewer than two bids per tender (so the
coefficient of variation, skewness, and kurtosis are undefined). Excluding
the direct CADE defendants by upstream label curation leaves an evaluation
pool of $16{,}731$ firms, of which $\valBidPosRetained$ are always-loser
cobidder positives. \textbf{All $\valBidPosRetained$ positive (cobidder)
labels are non-missing on every Imhof moment}: the support loss falls
entirely on the negatives, so there is no positive attrition. The seven
moments cover four families---dispersion, level/relative spread, shape, and
rank/order---and require the full within-tender bid record, which the award
layer does not contain. The full feature dictionary, units-of-observation
and per-feature missingness tables are in Online Supplement~E.

%--------------------------------------------------------------------
\subsection{Model Specification}
\label{app:bid_model_submission}

The learner is a random forest (\texttt{ranger}, $500$ trees,
\texttt{probability = TRUE}), with no hyperparameter tuning and no
standardization (the forest is scale-invariant); missingness is handled
listwise through the pool restriction. The outcome is binary
\texttt{is\_cade} ($1$ for an always-loser cobidder in the candidate pool,
not a direct defendant). Evaluation is by $5$-fold cross-validation, each
firm receiving one held-out predicted probability; AUC and DeLong tests are
computed on these out-of-fold firm-level predictions. Four designs of
increasing stringency are run---pooled-random (a diagnostic upper bound),
case-grouped (the honest robustness test), exclude-label-defining-tenders
(the contamination check), and environment holdout. Strict bid-timing
validation is \emph{blocked}: the within-tender moments cannot be cleanly
time-limited without re-deriving them on a per-period bid panel. The full
spec-by-spec audit and fold compositions are in Online Supplement~E.

%--------------------------------------------------------------------
\subsection{Leakage, Contamination, and Performance}
\label{app:bid_perf_submission}

The award score is \emph{clean} (a CADE-label-independent rank on
\texttt{log1p(tenders\_count)}). The bid and combined forests under
full-period features and random folds are \textbf{diagnostic-only}: random
folds are optimistic, and the firm-mean moments include label-defining
tenders. The contamination is small: excluding label-defining tenders
leaves the bid ROC essentially unchanged ($0.717 \to 0.713$), so the
fragility is case clustering, not feature--target overlap; the bid ROC
itself falls from $0.717$ (random) to $0.626$ (case-grouped), while the
combined model's case-grouped PR-AUC ($0.103$) falls below award-only
($0.143$). The headline ranking comparison---the cheap award layer's
comparable discrimination at lower cost, case-fragile complementarity, and
the moderate award/bid rank correlation---is reported and interpreted in
the main text (Section~6.1, Table~5). The full design-by-design
performance table, leakage audit, complementarity, and calibration tables
are in Online Supplement~E (bid-layer performance table).

%--------------------------------------------------------------------
\subsection{Limitations}
\label{app:bid_limitations_submission}

The bid-layer benchmark is a full-observability diagnostic, not a first
stage: the Imhof moments require all within-tender bids, which is exactly
the information cost the award-layer screen avoids; the headline
pooled-random numbers carry random-CV optimism; strict bid-timing
validation is \emph{blocked}, because within-tender dispersion moments
cannot be cleanly restricted to a pre-target window; the benchmark is
deliberately modest (untuned, unstandardized Imhof--Wallimann moments in a
random forest), not a state-of-the-art horse race; and the target is the
$\valBidPosRetained$ always-loser cobidders, an adjudication-anchored
exposure measure, so false positives are unlabeled firms and nothing here
constitutes proof of conduct.
